\section{Discussão e trabalho futuro}
\label{sec:discussao}

\subsection{O que este trabalho estabelece}

\begin{itemize}
  \item Um método de fronteira imersa com forçamento direto, verificado nos seus
        operadores de transferência independentemente do solver, em $np=1,2,3$.
  \item Que a formulação defasada tem erro $O(\dt/h)$ --- número de Courant ---
        e que refinar a malha sem corrigi-la \emph{piora} o não escorregamento.
  \item Que a rota pós-preditor com forçamento iterado reduz o resíduo em
        $41\times$ e permite $\dt$ vinte vezes maior para a mesma acurácia.
  \item $\Cd$ com erro de $+2{,}6\%$ a 20 células por diâmetro e $+0{,}4\%$ a 40,
        contra o valor de referência do \emph{benchmark}.
\end{itemize}

\subsection{O que este trabalho não estabelece}

\begin{itemize}
  \item \textbf{A inclinação de convergência.} Dois pontos não a determinam, e o
        valor obtido é quase o dobro do publicado. A terceira resolução está em
        curso.
  \item \textbf{Nada sobre $\Cl$.} O suporte do núcleo é maior que a assimetria
        que o produz.
  \item \textbf{Nada em 3D quantitativamente.} A geometria 3D-1Z existe e roda,
        mas a 10 células por diâmetro e sem comparação com referência.
  \item \textbf{Nada sobre corpo móvel ou deformável.} Ver abaixo.
  \item \textbf{O termo de inércia do fluido interno.} O $\Cd$ calculado ignora a
        inércia do fluido ``fantasma'' dentro do corpo, que Uhlmann contabiliza à
        parte. Para corpo fixo em regime permanente o termo se anula --- e o
        escoamento estacionou ---, de modo que o número é legítimo aqui; num
        transiente, não seria.
\end{itemize}

\subsection{Contorno rígido e interface entre fluidos são problemas distintos}

``Fronteira imersa'' cobre duas linhagens que compartilham a maquinaria e
divergem em quase tudo o mais:

\begin{center}
\begin{tabular}{lll}
  \toprule
   & \textbf{contorno rígido} & \textbf{interface entre fluidos} \\
  \midrule
  força        & multiplicador de Lagrange & constitutiva ($\sigma\kappa\mathbf{n}$) \\
  marcadores   & fixos                     & advectados \\
  conectividade& dispensável               & \textbf{indispensável} (curvatura) \\
  remalhamento & não                       & sim, ou a interface fica permeável \\
  \bottomrule
\end{tabular}
\end{center}

Os operadores de transferência servem aos dois sem diferença. A distribuição por
posse euleriana, porém, \textbf{destrói a ordem da curva} --- correto e barato
para o caso rígido, inviável para o outro, que precisa de vizinhos para calcular
curvatura.

\subsection{Refinamento local}

É o caminho da literatura: o erro é de ordem~1 \emph{junto} à interface e ordem~2
longe dela~\cite{lai2000,griffith2005}, e o método adaptativo de
Roma--Peskin--Berger~\cite{roma1999} --- fonte do núcleo aqui usado --- existe
exatamente para isso. Em 3D é a única rota viável: 40 células por diâmetro
uniformes custariam $\sim$27 milhões de células contra $\sim$420 mil hoje.

Três ressalvas, na ordem em que doem:

\begin{enumerate}
  \item A enumeração do suporte pressupõe um $h$ único e global. Numa malha
        graduada isso não é ajuste, é reescrita.
  \item O núcleo na fronteira de refinamento não preserva de graça a partição da
        unidade nem a propriedade adjunta --- que são precisamente as duas
        propriedades que os testes verificam.
  \item \textbf{Refinamento local reduz células, não passos.} O $\dt$ é global e
        ditado pela menor célula. O custo de refinar é $(1/h)^{D+1}$, e o
        refinamento local ataca apenas $D$ desse expoente.
\end{enumerate}

\subsection{Próximos passos}

\begin{enumerate}
  \item Terceira resolução (80 células/$D$), para decidir a inclinação de
        convergência. \emph{Em curso.}
  \item Se a inclinação se confirmar, refinamento local, com os pré-requisitos
        acima.
  \item Verificação 3D quantitativa contra o \emph{benchmark} 3D-1Z.
  \item Comparação com a rota de Lai--Peskin (ponto médio, duas avaliações de
        força por passo, sem iterar), que é a alternativa estrutural ao
        forçamento iterado.
\end{enumerate}
